\section{Discussion}
\label{sec:discussion}

\paragraph{Broader applicability.}
While our experiments use Reddit-style social threads, the three-pathway framework applies to any agentic system where the observe--generate--update loop creates opportunities for behavioral contamination.
Customer support agents that condition on conversation history, retrieval-augmented assistants that retrieve from potentially contaminated corpora, and tool-using agents that observe outputs from untrusted tools all share the same structural vulnerability.
The read/write control mechanism is model-agnostic and can be deployed as a wrapper around any agent runtime.

\paragraph{Limitations.}
Our study has several limitations.
First, the simulation uses LLMs role-playing as social media users, which may not perfectly capture how real users interact.
Second, Detoxify captures explicit toxicity but may miss subtler forms of harmful influence (coded language, strategic manipulation, gradual normalization); our LLM-based sentiment scorer partially addresses this but is itself imperfect.
Third, the two-model design (phenomenon on \texttt{gpt-4o-mini}, intervention on Llama-3.1-8B) means we cannot directly demonstrate DPO's effectiveness on the same model that exhibits the strongest phenomenon---though we validate that state-control interventions generalize across both models.
Fourth, we evaluate chains of at most 4 turns and trees of depth 2; longer interactions may exhibit different dynamics, including both stronger contamination (more state accumulation) and natural attenuation (topic drift).

\paragraph{Ethical considerations.}
We study toxic influence propagation to develop defenses, not to enable attacks.
Our simulations use synthetic agent-generated content with no real users involved.
The toxic system prompts used for $A_1$ are designed to elicit measurable effects within a controlled research setting and would not pass content moderation in production deployments.
All seed posts are drawn from publicly available Reddit data.


\section{Conclusion}
\label{sec:conclusion}

We have shown that toxicity in multi-agent social simulations is not a local output problem but a \emph{stateful propagation} phenomenon: a single toxic agent's influence persists through the evolving conversation state, contaminates compressed memory summaries, and shapes downstream agents' behavior even when those agents are neutral by design.
No single intervention---parameter unlearning, state sanitization, or memory control---suffices in isolation.
Parameter unlearning fails because toxic content re-enters through the state; state control fails because the model's learned biases regenerate toxic framing; memory sanitization fails because the raw transcript still carries contamination.
Our three-pathway framework addresses each channel simultaneously, achieving robust toxicity reduction under controlled evaluation with paired counterfactuals, dose--response analysis, and topology ablation.
These findings urge the community to treat agent safety as a \emph{systems problem} requiring coordinated intervention at the parameter, state, and memory levels---not merely a model alignment problem solvable by training alone.
